\section{The Information Problem and Preference Discovery} \label{sec:information}

\subsection{The Planner's Information Constraint}

A central difficulty in any system of economic planning is the
information problem emphasised in the socialist calculation debate
\citet{vonmises1920, hayek1945use}. In particular, the planner does
not observe the true utility function of households, nor the
underlying demand curves that govern consumption choices.

Let the true but unobserved utility function of households be

\[
W^*(\vec{C}) = W^*(C_1, \dots, C_n).
\]

The planner does not know the functional form of $W^*$, nor the
associated Marshallian demand functions $C_i^*(P, Y)$. Without this
information, the planner cannot directly solve the standard welfare
maximisation problem. Any implementable planning system must therefore
rely on a procedure for inferring consumer preferences from observable
behaviour.

\subsection{Ansatz Demand Specification}

To make the problem tractable, the planner adopts an ansatz for the
functional form of demand. Specifically, demand is approximated as
arising from a Constant Relative Risk Aversion (CRRA) utility
structure, in which each good $i$ carries a good-specific elasticity
parameter $\sigma_i > 0$, common across all households. The
resulting Marshallian demand for household $h$ takes the form:

\[
C_{h,i,t} = \left(\frac{w_h\,\alpha_{i,t,h}}{\mu_t\, P_{i,t}}\right)^{1/\sigma_i},
\]

where $\mu_t > 0$ is the marginal utility of income, common across
households and goods within period $t$, $\alpha_{i,t,h} > 0$ is a
time-varying preference weight for household $h$ over good $i$, and
$\sigma_i > 0$ is the good-specific relative risk aversion parameter.
Note that $\sigma_i$ is assumed identical across households but varies
freely across goods. This form is derived from utility maximisation of
a CRRA surrogate subject to the budget constraint
$\vec{P}_t^\top\vec{C}_{h,t} = Y_{h,t}$.

Aggregating across households, the aggregate demand for good $i$
satisfies:

\[
C_{i,t} = \sum_{h=1}^{n_h} C_{h,i,t}
         = \frac{1}{(\mu_t P_{i,t})^{1/\sigma_i}}
           \sum_{h=1}^{n_h} (w_h\,\alpha_{i,t,h})^{1/\sigma_i}.
\]

Defining the aggregate preference weight via the CES aggregation:
\[
\alpha^{agg}_{i,t} \;\equiv\; \left[\sum_{h=1}^{n_h}
  (w_h\,\alpha_{i,t,h})^{1/\sigma_i}\right]^{\!\sigma_i},
\]

aggregate demand takes the compact form:

\begin{equation}\label{eq:crra_demand}
C_{i,t} = \left(\frac{\alpha^{agg}_{i,t}}{\mu_t\, P_{i,t}}\right)^{1/\sigma_i}.
\end{equation}

The planner does not assert that true preferences are exactly CRRA;
rather, the weights $\alpha^{agg}_{i,t}$ are interpreted as a
parametric approximation to the unknown demand system, updated locally
in each period.

The key informational problem is therefore reduced to estimating the
vector

\[
\vec{\alpha}^{agg}_t = (\alpha^{agg}_{1,t},\dots,\alpha^{agg}_{n,t}),
\]

which evolves as the underlying preference structure changes. The
planner thus does not solve a full-information welfare maximisation
problem. Instead, the system implements a recursive adaptive control
process in which the working objective function is continuously
revised through parametric updating.

\subsection{Approximate Utility Maximisation}

Because the true utility function is unknown, the planner maximises a
CRRA surrogate using the current parameter estimate:

\[
{W}^{*}_t( \vec{w}_{h},\vec{C}_{h})
= \sum_{h = 1}^{n_{h}} w_{h}
  \left(\sum_{i=1}^n \alpha_{i,t,h}
        \,\frac{C_{h,i}^{1-\sigma_i}}{1-\sigma_i}\right)
\]

The resulting optimal consumption plan $\vec{C}^*_t$ is therefore
optimal only relative to the current approximation $\hat{W}_t$, not
necessarily relative to the true but unobserved preferences $W^*$.
This distinction is consequential: the plan is the best the planner
can do given available information, but it need not coincide with the
first-best social optimum. After each planning period, the realised
pattern of consumption provides new information about the true demand
structure, and the approximation is revised accordingly.

Here, we employ an approximate welfare function of the form:

\begin{equation}\label{eq:crra_welfare}
\hat{W}_t(\vec{C})
= \sum_{i=1}^n \alpha^{agg}_{i,t}
  \,\frac{C_{i}^{1-\sigma_i}}{1-\sigma_i}
\end{equation}

The marginal utility of aggregate consumption with respect to good $i$
is:
\[
\frac{\partial \hat{W}_t}{\partial C_i}
= \alpha^{agg}_{i,t}\, C_i^{-\sigma_i},
\]
and the marginal utility of income $\mu_t$ equates this to the shadow
price of good $i$:
\[
\mu_t \,P_{i,t} = \alpha^{agg}_{i,t}\, C_i^{-\sigma_i} = \pi_i,
\]
consistently with equation~\eqref{eq:crra_demand}.

\subsection{Parametric Preference Updating}

Within each planning cycle, prices adjust in response to excess demand
through the adjustment process described in Section~\ref{sec:dynamics}.
These price signals induce households to adjust their consumption
bundles, and observed expenditure patterns thereby contain information
about the underlying preference structure.

At the CRRA optimum, the first-order condition
$\alpha^{agg}_{i,t}\,C_i^{-\sigma_i} = \mu_t P_{i,t}$ implies:

\[
\frac{P_{i,t}\,(C_{i,t})^{\sigma_i}}{\displaystyle\sum_j P_{j,t}\,(C_{j,t})^{\sigma_j}}
= \frac{\alpha^{agg}_{i,t}}{\displaystyle\sum_j \alpha^{agg}_{j,t}}
= \alpha^{agg}_{i,t},
\]

where the last equality uses $\sum_i \alpha^{agg}_{i,t} = 1$.
Motivated by this relationship, the planner updates the preference
weights using observed marginal-utility-weighted expenditure shares:

\begin{equation} \label{eq:preference_update}
\alpha^{agg}_{i,t+1}
=
\frac{1}{3}\sum_{\tau = 1}^{3}
\frac{P_{t, \tau, i}\bigl(C^{m}_{\tau, t, i}\bigr)^{\sigma_i}}
     {\displaystyle\sum_j P_{t,\tau,j}\,\bigl(C^{m}_{\tau, t, j}\bigr)^{\sigma_j}}
\end{equation}

By construction this rule ensures $\sum_i \alpha_{i,t+1} = 1$ and
$\alpha_{i,t+1} > 0$ for all $i$, since all terms are non-negative and
the denominator is strictly positive. When $\sigma_i = 1$ for all
$i$, the rule reduces to ordinary expenditure shares, recovering the
Cobb--Douglas special case. Note that $\sigma_i$ enters the update
rule directly, so good-specific curvature is accounted for in the
revealed-preference inference.

It is important to note the limits of this updating rule. For a true
CRRA preference structure, the marginal-utility-weighted expenditure
shares are identically equal to $\alpha^{agg}_{i,t}$, so the rule
recovers the true parameters exactly. For any other preference
structure, the rule converges to the CRRA approximation that best fits
observed expenditure patterns, which will in general differ from the
true preference parameters. The persistent but bounded gap between
$\vec{\alpha}_t$ and the true preference vector is therefore a
structural feature of the approximation rather than a transient
estimation error, and is in fact the engine of the growth mechanism
described in Section~\ref{sec:growth}.

\subsection{A Cybernetic Preference Discovery Process}

The resulting system operates as a recursive preference approximation
mechanism. Each planning cycle performs three steps:

\begin{enumerate}
\item The planner computes an approximate optimal production plan
  using the current parameter estimate $\vec{\alpha}_t$.
\item Prices adjust through the t\^{a}tonnement process, generating
  observable consumption outcomes.
\item The parameter vector $\vec{\alpha}_t$ is updated from observed
  marginal-utility-weighted expenditure shares, revising the planner's
  approximation of consumer demand.
\end{enumerate}

Over time, this process functions as a cybernetic feedback system that
continuously adjusts the planning objective toward consistency with
observed household behaviour. Rather than requiring prior knowledge of
the utility function, the system refines its working approximation
adaptively. What the price signals reveal, in precise terms, is the
direction and magnitude of the gap between the current CRRA
approximation and the true demand structure --- information that is
then used to redirect investment toward sectors where demand is
systematically underestimated. This is a more limited claim than full
preference recovery, but it is sufficient for the adaptive planning
mechanism to function.